% =====================================================================
% STANDALONE DROP-IN SECTION FOR OVERLEAF
% =====================================================================
% Individual (per-instance) vehicle-indexed formulation for air assets,
% as implemented on the `individual-vehicle-indexing` branch
% (model/model.py, solve_stochastic_cvar with vehicle_formulation="individual",
% default flag configuration).
%
% This section is FULLY SELF-CONTAINED: it restates the sets, parameters,
% variables, objective, and all shared constraints so it can be read without
% flipping back to the aggregate formulation section. It is FORMULATION ONLY
% (no computational/tractability results).
%
% IT STANDS ALONGSIDE the existing experimental section (paper 3.12,
% "Individual Vehicle Indexing for Air Assets (Experimental)"). That earlier
% section describes a SUPERSEDED version of the model; this section reflects
% the current branch code. See the reconciliation remark at the end. Because
% both sections coexist, EVERY label here is prefixed `...:vif:...` to avoid
% clashing with the existing `eq:individualvehiclevar`, `eq:vehcons`,
% `eq:objective`, etc. Do a find/replace on the prefix if you prefer another.
%
% Requires: amsmath, amssymb (\mathbb, \geq, etc.), bbm or \mathbb for the
% indicator (uses \mathbb{1}); dsfont not required.
%
% NUMERIC PARAMETER VALUES (fleet sizes, speeds, degradation matrix, delta,
% beta, etc.) are unchanged from the main formulation's parameter/vehicle
% tables and are referenced, not duplicated, to avoid divergence between two
% copies. The one compact vehicle-parameter recap table below is included for
% readability; swap in \ref{}s to your existing tables if you prefer.
% =====================================================================

\section{Vehicle-Indexed Formulation for Air Assets (Complete)}
\label{sec:vif}

This section presents, in full and self-contained form, the extensive-form
two-stage stochastic program with \emph{individual (per-instance) indexing of
air-mode vehicles}. It supersedes the experimental sketch of
Section~\ref{sec:individualvehicles} and reflects the current implementation.
The aggregate formulation of Section~\ref{sec:formulation} remains the model
used for the primary results; the model below is an extension that resolves
the aggregate model's inability to distinguish one vehicle flying several legs
from several vehicles each flying one, at the cost of a substantially larger
integer program.

Individual indexing is scoped to the air mode only. Let
$K^{\mathrm{ind}} \subseteq K_{\mathrm{air}}$ denote the air-vehicle types
carrying per-instance identity; in the current configuration
$K^{\mathrm{ind}} = K_{\mathrm{air}} = \{\text{C-17}, \text{C-130J}\}$. Sea and
land types, and any air type not in $K^{\mathrm{ind}}$, retain the aggregate
integer-count representation unchanged. Vehicle identity is
\emph{scenario-local}: instance $\ell \in \{1,\dots,F_k\}$ denotes a vehicle
slot within a single scenario's realization, with no state carried across
scenarios, consistent with treating $\Omega$ as independent draws in which the
fleet begins each scenario at its base distribution.

\subsection{Sets}
\label{sec:vif:sets}

\begin{itemize}
  \item $N$: nodes, $|N| = 50$; $N^P \subseteq N$: PPL candidates, $|N^P| = 22$.
  \item $R = \{\text{food}, \text{water}\}$: resource types.
  \item $M = \{\text{sea}, \text{air}, \text{land}\}$: modes;
        $A_m \subseteq N \times N$: directed arcs on mode $m$.
  \item $T_i \subseteq M \times M$: feasible transfer mode-pairs at node $i$.
  \item $\Omega$: scenarios (Sample Average Approximation sample).
  \item $K_m$: vehicle types on mode $m$; $J_k \subseteq N^P$: basing-eligible
        nodes for type $k$.
  \item $K^{\mathrm{ind}} \subseteq K_{\mathrm{air}}$: air-vehicle types with
        per-instance indexing; for $k \in K^{\mathrm{ind}}$, instances are
        indexed $\ell \in \{1,\dots,F_k\}$.
\end{itemize}

\subsection{Parameters}
\label{sec:vif:params}

Demand $d^\omega_{ir}$; per-capita demand rate $\phi_r$; population $P_i$;
inventory ceiling $\bar{q}_{ir}$; safety-stock fraction $\rho$; availability
$a^\omega_{ir}\in\{0,1\}$; nominal and residual arc capacity $U_{m,ij,r}$,
$u^\omega_{m,ij,r}$; per-unit transport cost $c_{m,ij}$; arc distance
$\text{dist}_{ij}$; transfer capacity $\kappa_{i,m_1m_2}$ and cost multiplier
$\theta_{m_1m_2}$; unmet-demand penalty $\delta_{ir}$; site count cap
$P_{\max}$, activation cost $f_i$, selection budget $B$; CVaR confidence level
$\beta$; scenario weight $\pi^\omega = 1/|\Omega|$; vehicle-movement
tie-breaker $\varepsilon = 0.04$.

Vehicle parameters (Table~\ref{tab:vif:params}): fleet size $F_k$; effective
cruise speed $v_k$; distance budget $D_k = 3 v_k$ (the 72-hour window as three
days of range); turnaround penalty
$\psi_k = (\text{turnaround}_m / 24)\, v_k$; per-resource capacity
$\text{cap}_{k,r}$; initial basing count $b_{k,j}$ with
$\sum_{j} b_{k,j} = F_k$ and $b_{k,j} > 0 \Rightarrow j \in J_k \subseteq N^P$.

\paragraph{Instance-to-home assignment.}
For each individually-indexed type $k \in K^{\mathrm{ind}}$, the fleet's
per-node basing counts $b_{k,j}$ are expanded deterministically into
per-instance home assignments. Ordering the eligible base nodes and filling
instances in sequence, define the fixed home node
\begin{equation}
\zeta(k,\ell) \in J_k \subseteq N^P,
\qquad \ell \in \{1,\dots,F_k\},\ k \in K^{\mathrm{ind}},
\label{eq:vif:home}
\end{equation}
so that exactly $b_{k,j}$ instances have $\zeta(k,\ell) = j$. Because
$\sum_j b_{k,j} = F_k$, every instance receives exactly one home node, and
$\zeta(k,\ell)$ is a parameter (not a decision variable). This is the
per-instance refinement of the aggregate basing term $b_{k,j}\,p_j$.

\begin{table}[t]
\centering
\caption{Vehicle parameters (numeric values as in the main formulation's
vehicle table). Air types carry per-instance identity; sea/land are
aggregate.}
\label{tab:vif:params}
\begin{tabular}{llrrrr}
\hline
Type & Mode & $F_k$ & $v_k$ (km/day) & $D_k = 3v_k$ & $\psi_k$ \\
\hline
C-17     & air  & 12 & 11{,}662 & 34{,}986 & (2/24)$\,v_k$ \\
C-130J   & air  & 16 & 7{,}840  & 23{,}520 & (2/24)$\,v_k$ \\
LCU-1700 & sea  & 8  & 408      & 1{,}224  & (6/24)$\,v_k$ \\
M1083    & land & 60 & 1{,}116  & 3{,}348  & (1/24)$\,v_k$ \\
\hline
\end{tabular}
\end{table}

\subsection{Decision Variables}
\label{sec:vif:vars}

\paragraph{First stage (here-and-now).}
\begin{itemize}
  \item $p_i \in \{0,1\}$, $i \in N^P$: $1$ if node $i$ is selected as a PPL.
\end{itemize}

\paragraph{Second stage (recourse, per scenario $\omega$).}
\begin{itemize}
  \item $x^\omega_{m,ijr} \geq 0$: commodity-$r$ flow on mode-$m$ arc $(i,j)$.
  \item $\tau^\omega_{i,m_1m_2,r} \geq 0$: intermodal transfer of $r$ at node
        $i$ from mode $m_1$ to $m_2$.
  \item $z^\omega_{ir} \geq 0$: unmet demand of $r$ at node $i$.
  \item $y^\omega_{ir} \geq 0$: inventory of $r$ released from site $i$.
\end{itemize}

\paragraph{Vehicle variables.}
\begin{itemize}
  \item $n^\omega_{k,m,ij} \in \mathbb{Z}_{\geq 0}$: aggregate type-$k$ vehicle
        count on mode-$m$ arc $(i,j)$, for
        $k \in K_m,\ m \in \{\text{sea},\text{land}\}$ and for any air type
        $k \in K_{\mathrm{air}} \setminus K^{\mathrm{ind}}$. For
        $k \in K^{\mathrm{ind}}$ the symbol $n^\omega_{k,\text{air},ij}$ is
        retained but defined by the aggregation link
        \eqref{eq:vif:agglink} rather than optimized directly.
  \item $n^\omega_{k,\ell,\text{air},ij} \in \{0,1\}$: $1$ if instance $\ell$
        of air type $k \in K^{\mathrm{ind}}$ traverses air arc $(i,j)$ in
        scenario $\omega$.
\end{itemize}

\paragraph{Auxiliary variables.}
\begin{itemize}
  \item $\eta \in \mathbb{R}$: Value-at-Risk threshold; $\xi^\omega \geq 0$:
        CVaR excess loss; $L^\omega \geq 0$: scenario loss.
  \item $g^\omega_{k,j} \geq 0$: McCormick auxiliary linearizing the
        $p_j$-coupled turnaround exemption for aggregate types
        (\eqref{eq:vif:mccormick}).
  \item $g^\omega_{k,\ell} \geq 0$: McCormick auxiliary linearizing the
        home-node turnaround exemption for air instances
        (\eqref{eq:vif:mccormickind}).
\end{itemize}

\subsection{Objective and Scenario Loss}
\label{sec:vif:obj}

The scenario loss aggregates the unmet-demand penalty, transport cost across
all modal layers, intermodal transfer cost, and a small vehicle-movement
tie-breaker:
\begin{equation}
L^\omega =
\sum_{i \in N}\sum_{r \in R} \delta_{ir}\, z^\omega_{ir}
+ \sum_{m \in M}\sum_{(i,j) \in A_m}\sum_{r \in R} c_{m,ij}\, x^\omega_{m,ijr}
+ \sum_{i \in N}\sum_{(m_1,m_2) \in T_i}\sum_{r \in R}
      \theta_{m_1m_2}\, \tau^\omega_{i,m_1m_2,r}
+ \varepsilon \sum_{k}\sum_{m \in M}\sum_{(i,j) \in A_m} n^\omega_{k,m,ij}.
\label{eq:vif:loss}
\end{equation}
The final term ($\varepsilon = 0.04$) is a lexicographic tie-breaker,
negligible against $\delta_{ir}$ and without influence on delivery decisions;
it makes the solver strictly prefer fewer vehicle-arc traversals and no empty
repositioning legs when delivery outcomes are identical. For air types the
aggregate count $n^\omega_{k,\text{air},ij}$ appearing here equals
$\sum_\ell n^\omega_{k,\ell,\text{air},ij}$ by \eqref{eq:vif:agglink}, so the
tie-breaker counts individual instance movements consistently.

The model minimizes the Conditional Value-at-Risk of the loss at confidence
level $\beta$:
\begin{equation}
\min_{p,\,x,\,\tau,\,n,\,y,\,z,\,\eta,\,\xi,\,g}
\quad
\eta + \frac{1}{1-\beta} \sum_{\omega \in \Omega} \pi^\omega\, \xi^\omega .
\label{eq:vif:objective}
\end{equation}

\subsection{Constraints}
\label{sec:vif:constraints}

\paragraph{First-stage selection.}
\begin{align}
\sum_{i \in N^P} p_i &\leq P_{\max}, \label{eq:vif:pmax} \\
\sum_{i \in N^P} f_i\, p_i &\leq B, \label{eq:vif:budget} \\
p_i &\in \{0,1\} \quad \forall i \in N^P. \label{eq:vif:binary}
\end{align}
(The theater-wide inventory ceiling of the aggregate formulation is likewise
not imposed in the implementation and is omitted here; under $P_{\max}=3$ it
would remain non-binding.)

\paragraph{CVaR linearization (Rockafellar--Uryasev).}
\begin{align}
\xi^\omega &\geq L^\omega - \eta \quad \forall \omega \in \Omega,
  \label{eq:vif:cvar1} \\
\xi^\omega &\geq 0 \quad \forall \omega \in \Omega. \label{eq:vif:cvar2}
\end{align}

\paragraph{Inventory release.}
\begin{align}
y^\omega_{ir} &\leq (1-\rho)\, \bar{q}_{ir}\, a^\omega_{ir}\, p_i
  \quad \forall i \in N^P,\ r \in R,\ \omega \in \Omega, \label{eq:vif:release} \\
y^\omega_{ir} &= 0
  \quad \forall i \in N \setminus N^P,\ r \in R,\ \omega \in \Omega.
  \label{eq:vif:norelease}
\end{align}

\paragraph{Aggregate flow balance.}
\begin{align}
\sum_{m}\sum_{j:(j,i)\in A_m} x^\omega_{m,jir} + y^\omega_{ir} + z^\omega_{ir}
  &= d^\omega_{ir} + \sum_{m}\sum_{j:(i,j)\in A_m} x^\omega_{m,ijr}
  &&\forall i \in N^P,\ r,\ \omega, \label{eq:vif:balppl} \\
\sum_{m}\sum_{j:(j,i)\in A_m} x^\omega_{m,jir} + z^\omega_{ir}
  &= d^\omega_{ir} + \sum_{m}\sum_{j:(i,j)\in A_m} x^\omega_{m,ijr}
  &&\forall i \in N \setminus N^P,\ r,\ \omega. \label{eq:vif:balnon}
\end{align}

\paragraph{Per-mode outbound feasibility with transfer.}
\begin{align}
\sum_{j:(i,j)\in A_m} x^\omega_{m,ijr}
  + \sum_{m_2:(m,m_2)\in T_i} \tau^\omega_{i,mm_2,r}
\;\leq\;&
\sum_{j:(j,i)\in A_m} x^\omega_{m,jir}
  + \sum_{m_1:(m_1,m)\in T_i} \tau^\omega_{i,m_1m,r}
  + \mathbb{1}[i \in N^P]\, y^\omega_{ir} \notag \\
&\forall i \in N,\ m \in M,\ r \in R,\ \omega \in \Omega. \label{eq:vif:permode}
\end{align}

\paragraph{Transfer backing.}
\begin{equation}
\sum_{m_2:(m,m_2)\in T_i} \tau^\omega_{i,mm_2,r}
\;\leq\;
\sum_{j:(j,i)\in A_m} x^\omega_{m,jir}
  + \mathbb{1}[i \in N^P]\, y^\omega_{ir}
\quad \forall i,\ m,\ r,\ \omega.
\label{eq:vif:transferbacking}
\end{equation}

\paragraph{Capacity.}
\begin{align}
x^\omega_{m,ijr} &\leq u^\omega_{m,ij,r}
  \quad \forall m,\ (i,j) \in A_m,\ r,\ \omega, \label{eq:vif:arccap} \\
\sum_{r \in R} \tau^\omega_{i,m_1m_2,r} &\leq \kappa_{i,m_1m_2}
  \quad \forall i,\ (m_1,m_2) \in T_i,\ \omega. \label{eq:vif:transfercap}
\end{align}

\paragraph{Aggregate vehicle conservation (sea, land, non-indexed air).}
For every aggregate type --- $k \in K_m$ with $m \in \{\text{sea},\text{land}\}$,
and $k \in K_{\mathrm{air}} \setminus K^{\mathrm{ind}}$ --- a vehicle may depart
a node only if present there through activated basing or prior arrival:
\begin{equation}
\sum_{i:(i,j)\in A_m} n^\omega_{k,m,ij} + b_{k,j}\, p_j
\;\geq\;
\sum_{j':(j,j')\in A_m} n^\omega_{k,m,jj'}
\quad \forall j \in N,\ m,\ \omega.
\label{eq:vif:vehcons}
\end{equation}
No separate fleet-cardinality constraint is imposed: vehicles are injected
into the flow only at base nodes via $b_{k,j}\,p_j$, so
\eqref{eq:vif:vehcons} together with the distance budget
\eqref{eq:vif:distbudget} bounds fleet usage without capping traversals.

\paragraph{Air aggregation link.}
For each individually-indexed air type, the aggregate count is recovered as a
sum over instances, so the vehicle-capacity constraint
\eqref{eq:vif:vehcap} needs no modification:
\begin{equation}
n^\omega_{k,\text{air},ij} = \sum_{\ell=1}^{F_k} n^\omega_{k,\ell,\text{air},ij}
\quad \forall k \in K^{\mathrm{ind}},\ (i,j) \in A_{\text{air}},\ \omega.
\label{eq:vif:agglink}
\end{equation}

\paragraph{Per-instance conservation (air).}
For each air instance $\ell$, departures from a node are backed only by
arrivals of that same instance or, at its home node, by basing coupled to site
selection:
\begin{equation}
\sum_{i:(i,j)\in A_{\text{air}}} n^\omega_{k,\ell,\text{air},ij}
  + \mathbb{1}\!\left[\zeta(k,\ell)=j\right] p_j
\;\geq\;
\sum_{j':(j,j')\in A_{\text{air}}} n^\omega_{k,\ell,\text{air},jj'}
\quad \forall k \in K^{\mathrm{ind}},\ \ell,\ j \in N,\ \omega.
\label{eq:vif:indvehcons}
\end{equation}
This is the sole constraint coupling a specific air instance's availability to
the first-stage decision, through the term
$\mathbb{1}[\zeta(k,\ell)=j]\,p_j$.

\paragraph{Vehicle-capacity-constrained flow.}
Flow on an arc is bounded by the per-resource capacity of the vehicles
committed to it, using the aggregate count on every mode (with air's count
supplied by \eqref{eq:vif:agglink}):
\begin{equation}
x^\omega_{m,ijr} \;\leq\; \sum_{k \in K_m} \text{cap}_{k,r}\, n^\omega_{k,m,ij}
\quad \forall m,\ (i,j) \in A_m,\ r,\ \omega.
\label{eq:vif:vehcap}
\end{equation}

\paragraph{Fleet-wide distance budget (aggregate types).}
For each aggregate type $k$ with mode $m$, total distance flown plus turnaround
penalties on non-exempt outbound flow cannot exceed the fleet's aggregate
budget $F_k D_k$:
\begin{equation}
\sum_{(i,j)\in A_m} \text{dist}_{ij}\, n^\omega_{k,m,ij}
+ \psi_k \Bigg( \sum_{j \in N} \sum_{j':(j,j')\in A_m} n^\omega_{k,m,jj'}
  \;-\; \sum_{j \in J_k} g^\omega_{k,j} \Bigg)
\;\leq\; F_k\, D_k
\quad \forall k \notin K^{\mathrm{ind}},\ \omega.
\label{eq:vif:distbudget}
\end{equation}
The turnaround exemption applies only at base nodes that are actually selected;
because it depends on the product of outbound flow and the first-stage binary
$p_j$, it is linearized by the McCormick envelope
\begin{align}
g^\omega_{k,j} &\leq F_k\, p_j, \label{eq:vif:mccormick} \\
g^\omega_{k,j} &\leq \sum_{j':(j,j')\in A_m} n^\omega_{k,m,jj'}, \notag \\
g^\omega_{k,j} &\geq \sum_{j':(j,j')\in A_m} n^\omega_{k,m,jj'} - F_k\,(1 - p_j),
  \quad g^\omega_{k,j} \geq 0,
  \qquad \forall j \in J_k,\ \omega, \notag
\end{align}
which is exact because $p_j$ is binary.

\paragraph{Per-instance distance budget (air).}
Each air instance is held to a single-vehicle budget $D_k$ (not $F_k D_k$):
\begin{equation}
\sum_{(i,j)\in A_{\text{air}}} \text{dist}_{ij}\, n^\omega_{k,\ell,\text{air},ij}
+ \psi_k \Bigg( \sum_{j \in N} \sum_{j':(j,j')\in A_{\text{air}}}
      n^\omega_{k,\ell,\text{air},jj'}
  \;-\; g^\omega_{k,\ell} \Bigg)
\;\leq\; D_k
\quad \forall k \in K^{\mathrm{ind}},\ \ell,\ \omega.
\label{eq:vif:indvehdistbudget}
\end{equation}
Because an instance has a single fixed home node $\zeta(k,\ell)$, the
turnaround exemption has exactly one candidate node and its McCormick envelope
uses a constant bound of $1$:
\begin{align}
g^\omega_{k,\ell} &\leq p_{\zeta(k,\ell)}, \label{eq:vif:mccormickind} \\
g^\omega_{k,\ell} &\leq \sum_{j':(\zeta(k,\ell),j')\in A_{\text{air}}}
      n^\omega_{k,\ell,\text{air},\zeta(k,\ell)j'}, \notag \\
g^\omega_{k,\ell} &\geq \sum_{j':(\zeta(k,\ell),j')\in A_{\text{air}}}
      n^\omega_{k,\ell,\text{air},\zeta(k,\ell)j'} - \bigl(1 - p_{\zeta(k,\ell)}\bigr),
  \quad g^\omega_{k,\ell} \geq 0,
  \qquad \forall k \in K^{\mathrm{ind}},\ \ell,\ \omega. \notag
\end{align}

\paragraph{Rooted subtour elimination (air).}
Per-instance conservation \eqref{eq:vif:indvehcons} alone does not prevent an
instance's active arcs from forming a cycle disconnected from its home node: a
closed loop $j_1 \to j_2 \to \dots \to j_1$ touching no home node satisfies
conservation locally (one arrival and one departure at each node) while
consuming no basing credit. Such phantom cycles are excluded by the family of
rooted subtour-elimination constraints: for every instance, no node subset
excluding its home may contain a full internal cycle,
\begin{equation}
\sum_{\substack{(i,j)\in A_{\text{air}}\\ i \in S,\ j \in S}}
  n^\omega_{k,\ell,\text{air},ij}
\;\leq\; |S| - 1
\quad
\forall k \in K^{\mathrm{ind}},\ \ell,\ \omega,\
\forall S \subseteq N \setminus \{\zeta(k,\ell)\},\ |S| \geq 2.
\label{eq:vif:sec}
\end{equation}
A legitimate route from the home node traverses any such subset $S$ as an
acyclic path, using at most $|S|-1$ internal arcs, and is never cut; a
disconnected cycle within $S$ uses $|S|$ internal arcs and is removed. This
family is exponentially large and is not imposed a priori. It is separated
lazily within branch-and-cut (Dantzig--Fulkerson--Johnson): at each
integer-feasible incumbent, a routine computes, for every $(k,\ell,\omega)$,
the set of nodes reachable from $\zeta(k,\ell)$ along that incumbent's active
air arcs; each maximal group of active arcs touching only unreached nodes
identifies a violated subset $S$, whose constraint \eqref{eq:vif:sec} is added
as a lazy cut.

\paragraph{Instance symmetry breaking (air).}
Instances of the same air type sharing a home node are interchangeable
(identical $D_k$, $\psi_k$, $\text{cap}_{k,r}$, and home), differing only by
label. To remove the resulting permutation symmetry, instances within each
$(k, \text{home})$ group are ordered by total outbound arc count. Let
$H^\omega_{k,\ell} = \sum_{j\in N}\sum_{j':(j,j')\in A_{\text{air}}}
n^\omega_{k,\ell,\text{air},jj'}$ denote instance $\ell$'s total outbound air
flow; then, for consecutive instances $\ell, \ell+1$ with
$\zeta(k,\ell) = \zeta(k,\ell+1)$,
\begin{equation}
H^\omega_{k,\ell} \;\leq\; H^\omega_{k,\ell+1}
\quad \forall k \in K^{\mathrm{ind}},\ \omega .
\label{eq:vif:symbreak}
\end{equation}
Relabeling instances within an interchangeable group to satisfy this ordering
changes nothing else, so at least one optimum always satisfies it. The
ordering is scoped strictly within same-home groups: instances at different
home nodes are structurally distinguishable and are never constrained against
one another.

\paragraph{Nonnegativity and integrality.}
\begin{align}
z^\omega_{ir},\ y^\omega_{ir},\ x^\omega_{m,ijr},\
  \tau^\omega_{i,m_1m_2,r},\ g^\omega_{k,j},\ g^\omega_{k,\ell},\ \xi^\omega
  &\geq 0, \label{eq:vif:nonneg} \\
n^\omega_{k,m,ij} &\in \mathbb{Z}_{\geq 0}
  \quad (k \notin K^{\mathrm{ind}}), \label{eq:vif:int} \\
n^\omega_{k,\ell,\text{air},ij} &\in \{0,1\}
  \quad (k \in K^{\mathrm{ind}}), \label{eq:vif:bin} \\
\eta &\in \mathbb{R}. \label{eq:vif:etafree}
\end{align}

\subsection{Modeling Remarks}
\label{sec:vif:remarks}

\paragraph{Relationship to the experimental section.}
This formulation supersedes the sketch of
Section~\ref{sec:individualvehicles}. Two differences are substantive rather
than presentational. First, the mechanism preventing disconnected phantom
cycles is the lazily-separated subtour family \eqref{eq:vif:sec}, not a static
single-departure constraint. An alternative static formulation --- a per-node
departure-indicator binary with at most one active departure node per instance
--- is present in the implementation but disabled by default, because it cannot
simultaneously permit multi-leg routing and exclude disconnected cycles;
constraint \eqref{eq:vif:sec} does both. Second, the per-instance turnaround
exemption \eqref{eq:vif:mccormickind} is stated here with its correct McCormick
envelope, which the earlier section omitted.

\paragraph{Single-sortie-from-home.}
Because the exemption envelope \eqref{eq:vif:mccormickind} uses a constant
bound of $1$, when an instance's home node is selected the envelope forces its
home-node outbound flow to at most one arc; each air instance therefore
departs its home node at most once per scenario. This is a modeling
restriction of the present formulation (it excludes a within-scenario
return-refuel-relaunch from the home base), accepted here as consistent with
the single-response-window interpretation of a scenario.

\paragraph{Symmetry safety versus the aggregate model.}
The aggregate formulation of Section~\ref{sec:formulation} is symmetry-safe by
construction because same-type vehicles are never individuated. Individual
indexing reintroduces permutation symmetry, which \eqref{eq:vif:symbreak}
partially removes; residual symmetry among tied instances remains and is left
to the solver.
